\vspace{-2pt}
\section{Case Studies}
\label{sec:case_study}
\vspace{-4pt}


\subsection{Qualitative Results}
In this section, we present the results from various tasks implemented using InternAgent.

\subsubsection{Visual Examples of Various Tasks}


\begin{figure}[t]
    \centering
    \includegraphics[width=\linewidth]{Figures/showcase_autochem.pdf}
    \vspace{-10pt}
    \caption{Visual Examples of AutoRYP Task.}
    \label{fig:showcase_AutoRYP}
\end{figure}

\begin{figure}[t]
    \centering
    \includegraphics[width=\linewidth]{Figures/showcase_Molecule3D.pdf}
    \vspace{-10pt}
    \caption{Visual Examples of AutoMD Task.}
    \label{fig:showcase_AutoMD}
\end{figure}


We present showcases for three distinct tasks: AutoRYP, AutoMD, and AutoPower, to highlight the innovative methodologies discovered and their applications. These showcases are illustrated in Figs. ~\ref{fig:showcase_AutoRYP}, ~\ref{fig:showcase_AutoMD}, and ~\ref{fig:showcase_AutoPower}, respectively. Each task demonstrates a unique approach to solving complex problems, showcasing the potential impact of InternAgent across different scientific domains.

In the AutoRYP task illustrated in Fig.~\ref{fig:showcase_AutoRYP}, InternAgent autonomously discovered an innovative approach called "Adaptive Dual-Attention Graph-Transformer with Dynamic Freezing" for predicting chemical reaction yields. This method effectively integrates SMILES and graph-derived descriptors using a hybrid graph-transformer network, incorporating hierarchical attention mechanisms to enhance accuracy while minimizing overfitting. The approach features a Dual-Attention Fusion Mechanism (DAFM) that systematically combines token and graph embeddings with reaction conditions, ensuring effective information flow across different representations. Furthermore, a dynamic layer freezing mechanism, based on gradient magnitudes, optimizes which layers are trained, thereby enhancing generalization in low-data scenarios. Implemented with self-attention and cross-modality attention modules, this system not only combines 1D and 2D molecular representations effectively but also improves prediction accuracy and model adaptability. Consequently, this showcase underscores the method's potential for advancing research in complex chemical tasks using deep learning.

In the AutoMD task illustrated in Fig.~\ref{fig:showcase_AutoMD}, a novel framework called "Hierarchical Equivariant Directional Graph Encoder" (HEDGE-Net) has been autonomously discovered for predicting molecular energy and forces. This approach utilizes SE(3)-equivariant graph neural networks with hierarchical geometric self-attention and multi-hop message enrichment. By integrating angular and directional features into aggregated substructures, the method captures interacting atomic patterns and propagates dynamic weight updates, aligning with both local and global molecular geometries. The core of this method, the Geometry-Enhanced Directional Attention (GEDA) mechanism, ensures SE(3)-equivariance, enabling precise predictions for complex molecular systems at both atomic and substructural scales. Implemented with advanced message passing techniques, HEDGE-Net effectively combines directional and substructural information, enhancing scalability and precision in molecular modeling. This showcases the method's potential to advance research in complex molecular tasks using deep learning techniques.


\subsubsection{Visual Examples of Experimental Planning and Adaptive  Evolution}
\label{sec:case_exp_evo}

\begin{figure}[t]
    \centering
    \includegraphics[width=0.98\linewidth]{Figures/experimental_evolution_3dcls.pdf}
    \vspace{-8pt}
    \caption{Visual Examples of Experimental Planning and Adaptive Evolution on Auto3DCls task.}
    \label{fig:run5_showcase_3dcls}
\end{figure}

To further illustrate the practical utility of our experimental planning and adaptive evolution framework as described in Sec.~\ref{exp_evo}, we present some concrete examples of its application in the development and optimization for 3D point cloud classification and transcription prediction for perturbation response. Fig.~\ref{fig:run5_showcase_3dcls} and Fig.~\ref{fig:run5_showcase_tppr} visually summarize the stepwise experimental planning and adaptive evolution process that guided the implementation and refinement of our method.

\subsection{Human Evaluation}

Table~\ref{tab:idea_human_eval} compares the novelty of ideas generated by our InternAgent and AI-Scientist-V2~\citep{yamada2025ai}, across various research tasks. Each task involves generating 20 ideas, which are evaluated by five qualified reviewers. The assessments focus on four criteria: soundness, contribution, overall rating, and confidence. For each research task, the average scores of the 20 ideas are reported.

In the Reaction Yield Prediction task, InternAgent outperforms AI-Scientist-V2 in all aspects, especially in overall rating and soundness. Similarly, for 2D Semantic Segmentation, InternAgent shows better idea generation ability, particularly in soundness and overall rating. In 2D Image Classification and Point Cloud Autonomous Driving, InternAgent scores higher across all criteria, indicating a consistent advantage over AI-Scientist-V2 in generating more effective and novel ideas.


\begin{table}[t]
\centering
\renewcommand{\arraystretch}{1.2}
\caption{From the perspectives of soundness, contribution, and overall, we compare the novelty of ideas generated by InternAgent and AI-Scientist-V2~\citep{yamada2025ai}. For each research task, we generate 20 ideas. Each idea is scored by 5 qualified reviewers, and the final score for each task is reported as the average score of all 20 ideas. The detailed scores for each idea can be found in the Appendix~\ref{app:criteria}.}
\vspace{-6pt}
\label{tab:idea_human_eval}
\resizebox{\textwidth}{!}{
\begin{tabular}{l|l|cccc}
\toprule
\textbf{Research Task} & \textbf{Idea-gen Method} & \textbf{Soundness} & \textbf{Contribution} & \textbf{Overall} & \textbf{Confidence}  \\
\midrule
Reaction Yield Prediction & AI-Scientist-V2 & 1.42 & 1.45 & 3.50 & 3.50 \\
Reaction Yield Prediction & InternAgent & \textbf{3.09} & \textbf{2.66} & \textbf{4.35}  & \textbf{4.00} \\ \midrule
2D Semantic Segmentation & AI-Scientist-V2 & 1.84  & 2.07  & 2.95  &  \textbf{3.64}  \\
2D Semantic Segmentation & InternAgent &  \textbf{2.41}   &  \textbf{2.35}  &  \textbf{4.05}  &  3.48   \\  \midrule
2D Image Classification & AI-Scientist-V2 & 2.78   &  2.82   &  4.40   & \textbf{3.87} \\
2D Image Classification & InternAgent & \textbf{3.15}   &  \textbf{3.10}  & \textbf{5.85}   & 3.32 \\  \midrule
Point Cloud Autonomous Driving & AI-Scientist-V2 & 2.15 & 2.47 & 3.10 & 3.94 \\
Point Cloud Autonomous Driving & InternAgent & \textbf{2.75} & \textbf{2.95} & \textbf{5.10} & \textbf{4.10} \\ 

\bottomrule
\end{tabular}
}
\end{table}